\chapter{需求分析}\label{chap:req}

本章从用户角色与典型场景出发，给出 QuickBoard 的功能需求与非功能需求，并明确系统边界与约束。

\section{用户角色与使用场景}
系统主要面向无需注册的访客用户，典型场景包括：课堂讨论中的板书推导、线上会议中的流程图与方案讨论、远程结对时的草图沟通等。该类场景的共同点是：进入速度优先、协作过程对稳定性敏感、对复杂权限与长期项目管理需求较弱。

\section{功能需求}
\subsection{房间与入口}
\begin{itemize}
  \item \textbf{创建房间}：用户点击创建按钮后，系统生成房间号并进入白板页面。
  \item \textbf{加入房间}：用户可输入房间号或粘贴邀请链接进入对应房间；输入非法时应给出明确提示。
  \item \textbf{最近房间}：记录用户最近进入的房间列表，支持一键进入与移除记录。
\end{itemize}

\subsection{白板编辑}
\begin{itemize}
  \item \textbf{基础工具}：选择、画笔、形状（如矩形/圆形）、橡皮擦、文本等。
  \item \textbf{画布操作}：平移与缩放、对象选择与移动、删除与基本属性调整。
  \item \textbf{导出能力}：支持将白板内容导出为图片或数据文件，便于保存与分享。
\end{itemize}

\subsection{实时协作}
\begin{itemize}
  \item \textbf{多人同步}：同一房间的多个用户可看到彼此的绘制结果与对象变更。
  \item \textbf{重连与补包}：断线后能够自动重连，并尽量恢复到最新状态，避免刷新回滚。
\end{itemize}

\subsection{公式识别与编辑}
\begin{itemize}
  \item \textbf{公式识别}：用户框选区域后发送至 OCR 服务识别，返回 LaTeX。
  \item \textbf{确认插入}：识别结果需支持用户确认后再插入画布，避免误识别污染画布。
  \item \textbf{二次编辑}：插入后的公式可双击进入编辑，实时预览渲染效果并更新画布显示。
\end{itemize}

\section{典型用例}
为使需求具备可验证性，本课题将核心流程抽象为代表性用例，并在第 \ref{chap:eval} 章通过实验与场景测试进行验证。

\subsection{用例A：创建房间并分享}
\textbf{前置条件}：用户打开系统首页。\par
\textbf{基本流程}：点击创建按钮生成房间号并进入画板；页面展示可分享链接/房间号；其他用户通过链接或房间号加入同一房间。\par
\textbf{验收要点}：进入路径短、房间号易输入、分享后可直接进入同一协作域。

\subsection{用例B：多人实时绘制与撤销重做}
\textbf{前置条件}：两名或以上用户处于同一房间。\par
\textbf{基本流程}：用户 A 绘制路径/形状/文本并产生对象增量；服务端裁决后广播；用户 B 端出现一致对象；撤销/重做操作同步生效。\par
\textbf{验收要点}：对象集合一致、撤销重做行为多端一致。

\subsection{用例C：断线重连对齐}
\textbf{前置条件}：房间存在持续编辑且某客户端发生断线。\par
\textbf{基本流程}：客户端恢复连接后 join-room 携带 clientVersion；服务端根据 roomVersion 与增量日志能力发送 sync-delta 或 sync-state；客户端对齐后继续协作。\par
\textbf{验收要点}：断线窗口结束后最终收敛，且对齐路径明确。

\subsection{用例D：公式识别与插入}
\textbf{前置条件}：用户在画布上写出公式或导入手写内容。\par
\textbf{基本流程}：框选区域并发起识别；返回 LaTeX 或错误码；用户确认/修改后插入为公式对象；远端同步显示并可再次编辑。\par
\textbf{验收要点}：成功路径闭环；失败路径提示明确；可通过 traceId 与可选 debug 图定位失败原因。

\section{非功能需求}
\subsection{可用性与易用性}
系统应保证首屏聚焦主任务（创建/加入房间），并在操作反馈上保持清晰。关键流程应尽量减少弹窗与中断式交互，错误提示应直观且可恢复。

\subsection{实时性与一致性}
在网络条件正常时，多人协作的更新应具有较低延迟；在并发冲突与消息乱序场景下，系统需提供可解释的裁决策略，并尽量避免状态回滚。

\subsection{稳定性与性能}
系统需在常见浏览器环境下稳定运行。白板渲染应在多对象场景下保持可交互性，公式识别在失败时应及时返回并允许用户继续操作。

\subsection{可验证指标}
为提高论文论证的可复核性，本课题将部分非功能需求转化为可量化指标：
\begin{itemize}
  \item \textbf{最终收敛}：在丢包与断线场景下，多端结束时的权威对象集合应收敛到单一版本（结束时哈希种类数为 1）。
  \item \textbf{渲染性能}：在 5000 对象规模的基准场景下，启用视口裁剪与对象缓存后平均渲染耗时应明显下降。
  \item \textbf{过程降载}：在持续书写输入下，过程预览链路应显著降低消息数量，并保证结束信号正确。
\end{itemize}

\begin{table}[htbp]
  \centering
  \caption{核心需求与验证方式}\label{tab:req-verify}
  \renewcommand{\arraystretch}{1.35}
  \begin{tabular}{lll}
    \toprule
    \textbf{需求项} & \textbf{验证方式} & \textbf{证据来源} \\
    \midrule
    弱网最终收敛 & 一致性对照实验 & \texttt{doc/test\_reports/quickboard\_sync\_benchmark} \\
    大对象渲染可用 & 渲染性能基准 & \texttt{doc/test\_reports/quickboard\_render\_perf} \\
    过程预览不拥塞 & Ghost Brush 降载统计 & \texttt{doc/test\_reports/quickboard\_ghost\_brush} \\
    OCR 可用且可定位 & 场景测试 + 错误码/traceId & 项目日志与截图 \\
    \bottomrule
  \end{tabular}
\end{table}

\subsection{安全性与隐私}
系统默认无需用户身份信息，房间链接作为共享入口应避免泄露敏感数据；对于本地保存的最近房间记录，应明确其仅存储在浏览器本地存储中。

\section{约束与边界}
本项目聚焦本科毕业设计的实现与验证，不包含账号体系、复杂权限管理与长期存储的团队空间等平台化能力。协作一致性以“工程可用与体验稳定”为优先目标，不追求强一致性下的复杂冲突合并算法。
